\documentclass[12pt]{article}
\usepackage[margin=0.6in,top=1.15in,bottom=0.55in]{geometry}
\usepackage{lmodern}
\usepackage{microtype}
\usepackage{titlesec}
\usepackage{enumitem}
\usepackage[hidelinks]{hyperref}
\usepackage{../../latex/grader-worksheet}

\setlength{\parindent}{0pt}
\setlength{\parskip}{0.25em}
\titleformat{\section}{\large\bfseries\scshape}{}{0pt}{}

\GraderSetup{assignment-id=U1L11LE,template-revision=1}

\begin{document}

\begin{center}
{\LARGE\bfseries Week 11 Lit Example Worksheet}\par\vspace{0.05em}
{\large\scshape Testing Shakespeare's Comparisons}\par\vspace{0.12em}
\rule{0.78\textwidth}{0.5pt}
\end{center}

\textbf{Purpose:} Apply Week 11's method to Brutus and Macbeth. Decide what the comparisons prove and where they overextend.

\textbf{Directions:} Use the Lit Example Reader. Do not retell the plot; test the comparison as proof.

\section*{Part 1: Brutus's Source and Target}

\textbf{Question 1.1}\\[-0.15em]
In Brutus's serpent's egg comparison, what is the source case and what is the target case?

\GraderAnswerBox{Part 1 Q1}{1.95in}

\textbf{Question 1.2}\\[-0.15em]
What conclusion does Brutus want this comparison to support?

\GraderAnswerBox{Part 1 Q2}{2.05in}

\newpage

\section*{Part 2: Relevant Similarity and Disanalogy}

\textbf{Question 2.1}\\[-0.15em]
What relevant similarity makes Brutus's comparison persuasive?

\GraderAnswerBox{Part 2 Q1}{2.25in}

\textbf{Question 2.2}\\[-0.15em]
What important disanalogy limits Brutus's comparison?

\GraderAnswerBox{Part 2 Q2}{2.25in}

\newpage

\section*{Part 3: Does Brutus Prove Too Much?}

\textbf{Question 3.1}\\[-0.15em]
Does the serpent's egg analogy prove that Caesar must be killed, support only a narrower warning, or merely decorate Brutus's fear? Explain.

\GraderAnswerBox{Part 3 Q1}{2.55in}

\textbf{Question 3.2}\\[-0.15em]
Rewrite Brutus's conclusion so it matches what the analogy can honestly support, unless additional evidence is supplied.

\GraderAnswerBox{Part 3 Q2}{2.35in}

\newpage

\section*{Part 4: Macbeth's Bear-Like Comparison}

\textbf{Question 4.1}\\[-0.15em]
In Macbeth's bear-like comparison, identify the source case, target case, relevant similarity, and one disanalogy.

\GraderAnswerBox{Part 4 Q1}{2.45in}

\textbf{Question 4.2}\\[-0.15em]
Does Macbeth's comparison prove that he has no moral choice left, or does it mainly describe his sense of being trapped? Explain.

\GraderAnswerBox{Part 4 Q2}{2.45in}

\end{document}
